% ============ 03. Veredicto ============
\section{Veredicto global}

\begin{tcolorbox}[hallazgo]
\textbf{NO LISTO para entregar hoy como anteproyecto — LISTO para producir la
entrega con calidad alta en 3--5 días dirigidos.} Los tres evaluadores coinciden:
el equipo posee \textbf{materia prima de nivel Competente--Ejemplar} (cifras
propias verificadas, estado del arte validado con 0 referencias fabricadas,
conocimiento profundo de los datos), pero el \textbf{anteproyecto como documento
de 6--10 páginas en la estructura exigida no existe todavía}, no hay
\textbf{declaración de IA y reparto} (S3), y la \textbf{preparación individual
para la sustentación} (60\,\% del peso) no está verificada. Los dos criterios
\textbf{bloqueantes (D1 y D3)} están al filo de \emph{Insuficiente} no por falta
de evidencia, sino por falta de \textbf{redacción} del problema y de objetivos
nuevos.
\end{tcolorbox}

\begin{table}[htbp]
\centering
\caption{Resumen del bloque Documento (40\,\%) — niveles consensuados}
\label{tab:resumen}
\small
\begin{tabular}{lcccc}
\toprule
\textbf{Criterio} & \textbf{Peso} & \textbf{Nivel} & \textbf{Logro} & \textbf{Nota parcial (×5)} \\
\midrule
D1 Problema y contexto (bloqueante) & 12\,\% & En desarrollo & 0,6 & 0,36 \\
D2 Estado del arte & 9\,\% & En desarrollo & 0,6 & 0,27 \\
D3 Objetivos (bloqueante) & 8\,\% & En desarrollo & 0,6 & 0,24 \\
D4 Datos y método preliminar & 6\,\% & \textbf{Competente} & 0,8 & 0,24 \\
D5 Alcance, viabilidad y plan & 5\,\% & En desarrollo & 0,6 & 0,15 \\
\midrule
\multicolumn{4}{l}{\textbf{Total bloque Documento}} & \textbf{1,26 / 5,0} \\
\bottomrule
\end{tabular}
\end{table}

\begin{tcolorbox}[nota]
\textbf{Lectura:} 1,26/5,0 en el bloque de documento equivale al 63\,\% del
máximo del bloque (2,0 sobre la nota total). No incluye la sustentación
(60\,\%). El nivel ``En desarrollo'' refleja que \emph{hoy no existe el texto
pedido}, no la calidad de la evidencia subyacente: con la redacción correcta,
los mismos materiales apuntan a Competente--Ejemplar en D1, D2, D3 y D5.
\end{tcolorbox}
